\begin{ejercicio}
    \begin{enumerate}
    \item[(a)] Hagan un bosquejo de la región de integración de la integral a continuación e incluyan la mayor cantidad de información posible.
    \begin{equation}
        \int_{0}^{6} \int_{x/3}^{2} x\sqrt{y^3 + 1}\, dy\, dx
    \end{equation}
    \item[(b)] Describan verbalmente cómo se resolvería la integral en el orden indicado en el inciso (a) y luego cómo se resolvería invirtiendo el orden de integración. ¿En qué radica la diferencia?
\end{enumerate}
\end{ejercicio}

\begin{respuesta}
     \textbf{Integrantes involucrados: }
\end{respuesta}